% Reúne las reiteraciones del resumen extenso anterior; no sustituye las pruebas.
\subsection*{Mapa de resultados y dependencias}
El argumento contiene cinco niveles que conviene distinguir al leer. La
construcción aritmética proporciona los modos; sus lectores definen la forma
completa; las estimaciones establecen positividad en dominios precisos; el
criterio de Weil determina qué identificación adicional permite concluir RH.
El orden de anuncio hace visible el resultado principal, mientras el cuerpo
conserva el orden de sus dependencias.

\begin{center}
\renewcommand{\arraystretch}{1.22}
\begin{tabular}{@{}p{0.21\linewidth}p{0.48\linewidth}p{0.23\linewidth}@{}}
\hline
\textbf{Nivel} & \textbf{Resultado y función en el argumento} & \textbf{Residencia}\\
\hline
Construcción común &
Historias, hojas, memoria e incidencias permanecen en un mismo sistema de
refinamiento. Centro, frontera levantada y vacancias especifican qué
información transmite cada lectura. &
Secciones~\ref{lectura:manuscrito-05} y~\ref{sec:centro-frontera}.\\
Aritmética y funcional &
Producto nativo, formas normales e irreducibles; ocupaciones, producto de
Euler, momentos primo--potencia, partes finitas y continuación theta--Mellin. &
Secciones~\ref{lectura:manuscrito-01}--\ref{lectura:manuscrito-04};
inicio de~\ref{lectura:manuscrito-03}.\\
Reconstrucción y transporte &
Las columnas conservan conjuntamente los términos primo, gamma y polar.
La cola gamma reconstruye la prueba. La compresión, la memoria y Cayley
conservan sus balances y productos cruzados. &
Secciones~\ref{lectura:manuscrito-03},
\ref{sec:transporte-momentos} y~\ref{sec:lectores-gamma}.\\
Signo de la forma &
Coercividad del ensamblaje especificado, control de detalles y reducción de
Schur. La ventana completa de longitud \(1/9\) satisface
\(\mathscr W(f,f)\geq\frac12\|f\|_2^2\). &
Sección~\ref{lectura:manuscrito-03} y
teoremas~\ref{schur-add-schur-exacto},
\ref{schur-add:coercividad-ventana}.\\
Criterio global &
La identificación de los momentos completos, con todos sus cruces y cobertura
del dominio de prueba, permitiría aplicar el criterio de Weil. No se sustituye
por las cotas locales anteriores. &
Sección~\ref{lectura:manuscrito-03};
subsecciones~\ref{tm-add-c5} y~\ref{tm-add-m7}.\\
\hline
\end{tabular}
\end{center}

Dos distinciones acompañan este mapa. En aritmética, el período de una acción
residual no sustituye la irreducibilidad del producto. En análisis, una
factorización positiva de los momentos de memoria no se identifica
automáticamente con la de los momentos aritméticos. Las constantes
\(8/81\) y \(5/9\) pertenecen, respectivamente, al balance de compresión y
a la razón modal de memoria; sus pruebas conservan esas funciones distintas.

El centro \((5,5)\), la firma regional \(555555\), sus tres hojas de
continuación y los calendarios \(21/22\) y \(6/7\) se explican en su posición
constructiva. Las lecturas \(72,27,77,22\), la frontera levantada y la
completación \(729+271\) no se agrupan por coincidencia numérica: cada una
conserva la operación y el dominio que la produce. Su función en este
artículo es precisar la información transportable antes de pasar al
problema espectral.

La composición de lectores conserva además dos operaciones distintas:
el producto diagonal de tablas residuales y el producto tensorial de
sus momentos, que produce convolución de divisores. El transporte de
Stieltjes incluye la contribución afín del polo. La jerarquía de
Bendersky--Glaisher relaciona las derivadas regularizadas en niveles
negativos pares con los valores positivos impares del mismo funcional.
Los encierros por corte de irreducibles complementan las cotas por
profundidad, con sus errores respectivos explícitos
(secciones~\ref{ix-add:corte-primos},
\ref{ix-add:composicion-lectores} y~\ref{ix-add:bendersky}).

Las definiciones y pruebas utilizadas se reúnen en las secciones indicadas
y en los apéndices A--C, con los dominios de cada resultado explícitos.
Las pruebas locales, los resultados funcionales y los enunciados
condicionales mantienen aquí su alcance propio.
